\section{Formal set theory}\label{sec:formal_set_theory}

We start with some basics of set theory from the perspective of \fullref{ch:mathematical_logic}. We will work in a less formal metalanguage afterwards.

\paragraph{Axiomatic set theory}

\begin{definition}\label{def:axiomatic_set_theory}\mimprovised
  \term[en=axiomatic set theory (\cite[xii]{Hinman2005Logic})]{Axiomatic set theory} may refer to any \hyperref[def:logical_theory]{logical theory} based on a single binary \hyperref[con:predicate_logic/predicate]{predicate} \( \synin \), called \term[en=membership relation (\cite[5]{Jech2006SetTheory})]{set membership}\fnote{Other operations and predicates can be added as \hyperref[con:syntactic_abbreviation]{abbreviations} or \hyperref[def:fol_definitional_extension]{definitional extensions} --- see \cref{def:set_theory_definition} --- but we must be careful not to add additional structure.}. The adjective \enquote{axiomatic} is meant to highlight that the theory is stated in a fixed \hyperref[con:logical_system]{logical system}, as opposed to the informal theory like the one Cantor originally described in \cite{Кантор1985Многообразия}.

  The theories we will consider will be first-order adaptations, in the sense of \cref{def:fol_adaptation_of_monadic_sol_theory}, of \hyperref[def:second_order_logic]{monadic second-order} \hyperref[def:hol_theory]{theories}. This particular choice of formalisms is due to established practice --- see \cref{rem:first_order_zfc}.

  \begin{thmenum}
    \thmitem{def:axiomatic_set_theory/extensionality}\mcite[16]{Kunen2013SetTheory} Different set theories differ by their axioms --- these are existential axioms that postulate that a particular construction we can informally describe is actually well-defined.

    An exception is the \term[ru=аксиома объёмности (\cite[180]{Герасимов2014Вычислимость})]{axiom of extensionality}, which expresses the essence of a set --- namely, that a set is completely determined by its members\fnote{Technically, we must take the \hyperref[def:fol_quantifier_closure]{universal closure} of \eqref{eq:def:axiomatic_set_theory/extensionality}.}:
    \begin{equation}\label{eq:def:axiomatic_set_theory/extensionality}
      \qforall \synz (\synz \synin \synx \syniff \synz \synin \syny) \synimplies \synx \syneq \syny.
    \end{equation}

    The converse follows by a simple application of \ref{inf:def:fol_natural_deduction/equality/elim}.

    \thmitem{def:axiomatic_set_theory/atoms}\mimprovised In \cite{Zermelo2010Collection:GrundlagenDerMengenlehre}, where Zermelo introduces his axiomatization of set theory, he assumes the existence of a \enquote{Bereich} (\enquote{domain}) consisting of \enquote{Dinge} (\enquote{things}). Some of these things are assumed to be \enquote{Mengen} (\enquote{sets}).

    Later on, as set theory develops, the perspective shifts and objects that are not sets became secondary. They become known as \term[en=atom (\cite[27]{Schechter1997AnalysisHandbook})]{atoms} or \term[en=urelement (\cite[238]{Shoenfield1967MathematicalLogic})]{urelements}\fnote{The word \enquote{urelement} uses the German prefix \enquote{ur}, meaning \enquote{primordial}. According to \cite{MacTutor:EarliestUsesOfMathWordsU}, the term is borrowed from Zermelo's \cite{Zermelo2010Collection:GrenzzahlenUndMengenbereiche}.}

    Atoms can be formalized using \hyperref[def:many_sorted_fol]{many-sorted first-order logic} with a dedicated sort \( \syn\sigma \) (\enquote{\textel{σύνολο}}) for sets and a separate relation \( \synin_S: S \synimplies \syn\sigma \synimplies \syn\omicron \) for every sort \( S \).

    Since atoms are not essential for set theory, and since many-sorted logic introduces considerable overhead, it is customary to study unisorted set theory, where the only individuals are sets. We call such sets \term[en=pure (sets) (\cite[238]{Shoenfield1967MathematicalLogic})]{pure}.
  \end{thmenum}
\end{definition}
\begin{comments}
  \item See \cref{def:set_theory_definition} for a convenient syntax for describing sets.
  \item We discuss set existence in \cref{rem:set_existence}.
\end{comments}

\begin{remark}\label{rem:epsilon_and_set_membership}
  In \cite{MacTutor:EarliestUsesOfLogicSymbols}, the symbol \enquote{\( \in \)} for \hyperref[def:axiomatic_set_theory]{set membership} is attributed to Peano's \cite{Peano1889ArithmeticesPrincipia}. Peano himself claims to have derived the symbol from the Latin word \enquote{\textla{est}} (\enquote{it is}).

  Peano himself used the stylized epsilon symbol \enquote{\( \in \)}, however some later books like \cite{Hausdorff1935Mengenlehre} use \enquote{\( \varepsilon \)}. The usual epsilon occurs as late as half a century later, with \cite{Kelley1975GeneralTopology} and \cite{Halmos1976MeasureTheory}.

  The use of epsilon explains the name \fullref{thm:epsilon_induction}.
\end{remark}

\begin{remark}\label{rem:set_theory_acronyms}
  \hyperref[def:axiomatic_set_theory]{Axiomatic set theories} have several established acronyms:
  \begin{thmenum}
    \thmitem{rem:set_theory_acronyms/zfc} The ubiquitous acronym \logic{ZFC} is used for \hyperref[def:zfc]{Zermelo-Fraenkel set theory} with the axiom of choice. We list some variations.
    \begin{itemize}
      \thmitem{rem:set_theory_acronyms/zf} \logic{ZF} denotes \logic{ZFC} without the axiom of choice. The two theories are often introduced simultaneously, for example in
      \cite[510]{Hinman2005Logic},
      \cite[3]{Jech2006SetTheory},
      \cite[17]{Kunen2013SetTheory},
      \cite[253]{Shoenfield1967MathematicalLogic} and
      \cite[161]{КолмогоровДрагалин2006ЛогикаДополнительныеГлавы}.

      \thmitem{rem:set_theory_acronyms/zfcu} As discussed in \cref{def:grothendieck_universe}, for the purposes of \hyperref[ch:category_theory]{category theory}, we assume the existence of a \hyperref[def:grothendieck_universe]{Grothendieck universe} in addition to the axioms of \logic{ZFC}. We denote this theory by \hyperref[def:zfcu]{\logic{ZFC+U}}, although this is not an established notation.

      \Cref{thm:grothendieck_universe_iff_regular_strong_limit} implies that the Grothendieck universes are precisely the \hyperref[rem:strongly_inaccessible_cardinal]{strongly inaccessible cardinals}, which leads to the acronym \logic{ZFC+I} used in \cite[167]{Jech2006SetTheory} and \cite[16]{Shulman2008SetTheoryForCategoryTheory}.

      \thmitem{rem:set_theory_acronyms/zfcs} A variant specifically constructed for \hyperref[ch:category_theory]{category theory} is \logic{ZFC/S} (\enquote{\logic{ZFC} with smallness}), which we briefly discuss in \cref{rem:small_category/reflection}.
    \end{itemize}

    \thmitem{rem:set_theory_acronyms/nbg} \hyperref[def:nbg]{von Neumann-Bernays-G\"odel set theory} can be by different acronyms. For example, \logic{BG} is used in \cite[70]{Jech2006SetTheory}, \logic{NBG} is used in \cite[232]{Mendelson2015Logic}, while \logic{NGB}is used in \cite[151]{Kunen2013SetTheory}.

    In a footnote, \textcite[231]{Mendelson2015Logic} writes the following:
    \begin{displayquote}
      I coined the name NBG in honor of von Neumann, Bernays, and G\"odel. Paul Halmos, who favored the Zermelo-Fraenkel system, suggested that \enquote{NBG} stood for \enquote{No Bloody Good.}
    \end{displayquote}

    \thmitem{rem:set_theory_acronyms/mk} Another set theory is \logic{MK}, named after Morse and Kelley, and attributed by \textcite[\S 4.6.1]{Mendelson2015Logic} to the appendix of Kelley's \cite[201]{Kelley1975GeneralTopology}.

    We will not examine \logic{MK} in this monograph.
  \end{thmenum}
\end{remark}

\begin{definition}\label{def:set_theory_model}\mimprovised
  Let \( T \) be a (first-order) \hyperref[def:axiomatic_set_theory]{axiomatic set theory}, an \hyperref[def:fol_adaptation_of_monadic_sol_theory]{adaptation} of the second-order theory \( T_2 \).

  A \hyperref[def:fol_semantics/model]{first-order model} \( \mscrM = (M, I) \) of \( T \) is characterized mainly by the universe \( M \) and the membership interpretation \( I(\synin) \). Unless we are dealing with additional structure like \hyperref[def:axiomatic_set_theory/atoms]{urelements} or \hyperref[def:nbg]{classes}, both first-order and second-order models can be regarded as pairs \( (M, I(\synin)) \).

  \begin{thmenum}
    \thmitem{def:set_theory_model/standard} \Cref{def:fol_adaptation_of_monadic_sol_theory} suggests calling \( \mscrM \) \enquote{standard} if it is a model of \( T_2 \).

    This is well-suited for some theories like \hyperref[def:peano_arithmetic]{Peano arithmetic}, whose second-order variants are \hyperref[def:hol_theory/categorical]{categorical} --- all standard models are isomorphic, so for many purposes it is immaterial which one we work with.

    Unfortunately, as a set theory, \( T_2 \) is likely not \hyperref[def:hol_theory/categorical]{categorical}. For example, \fullref{thm:zermelos_quasi_categoricity_theorem} characterizes the models of \hyperref[def:zfc]{\( \logic{ZFC}_2 \)} as the \hyperref[def:grothendieck_universe]{Grothendieck universes}, of which there is a hierarchy. Some categorical second-order set theories are explored in \cite{HamkinsSolberg2020SetTheoreticReflection}.

    We refrain from speaking about \enquote{standard models} of set theory.

    \thmitem{def:set_theory_model/transitive}\mcite[def. I.15.1]{Kunen2013SetTheory} It is most natural to study models \( (M, {\in}) \), where \( M \) is a set and \( {\in} \) is the metatheoretic membership relation restricted to \( M \)\fnote{Kunen calls such models \enquote{\( \in \)-models}.}.

    This is by itself a weak condition because of the following subtle problem: \( B \in A \in M \) does not by itself imply that \( B \in M \). It is thus possible for metatheoretic members of \( A \) to not be considered members \emph{inside} \( (M, {\in}) \).

    If this cannot happen, we call the model \term[ru=транзитивная (модель) (\cite[\S 4.7.2]{КусраевКутателадзе2005БулевозначныйАнализ})]{transitive}. The universe of a transitive model is a transitive set in the sense of \cref{def:transitive_set}.
  \end{thmenum}
\end{definition}

\begin{definition}\label{def:set}\mimprovised
  We will refer to any element of any model of any \hyperref[def:axiomatic_set_theory]{axiomatic set theory} as a \term{set}.
\end{definition}
\begin{comments}
  \item Unless explicitly mentioned otherwise, sets are assumed to belong to a fixed \hyperref[def:set_theory_model/transitive]{transitive} model of \hyperref[def:zfcu]{\logic{ZFC+U}}.

  \item Because of the \hyperref[def:axiomatic_set_theory/extensionality]{axiom of extensionality}, a set is determined by its members, leading to the ubiquitous set-builder notation we discuss in \cref{def:set_theory_definition}. To show that a hypothetical set is well-defined, we only need to prove that a set with these exact members exists. See \cref{rem:set_existence}.
\end{comments}

\paragraph{Naive set theory}

\begin{remark}\label{rem:naive_set_theory}
  The phrase \enquote{naive set theory} can have two meanings:
  \begin{thmenum}
    \thmitem{rem:naive_set_theory/informal} \enquote{Naive} may refer to the informal (i.e. non-\hyperref[def:axiomatic_set_theory]{axiomatic}) set theory initially studied by Cantor. In this sense, the phrase is used, for example, in
    \cite[11]{Enderton1977SetTheory} and
    \cite[\S I.1]{Aluffi2009Algebra}.

    Paul Halmos describes his book \cite{Halmos1960NaiveSetTheory} as \enquote{axiomatic set theory from the naive point of view} because he explains \hyperref[def:zfc]{\logic{ZFC}} without focusing on formalities.

    \thmitem{rem:naive_set_theory/formal} \enquote{Naive set theory} may refer to the specific axiomatic set theory with unrestricted comprehension --- the one we state in \cref{def:naive_set_theory}. In this sense, the phrase is used in
    \cite[\S I.1]{КолмогоровДрагалин2006ЛогикаДополнительныеГлавы} and
    \cite[\S 3.4]{Герасимов2014Вычислимость}.

    \Textcite[\S 5.3.1]{Mimram2020ProgramEqualsProof} uses the phrase \enquote{naive set theory} for the axiomatic theory with unrestricted comprehension but without extensionality.

    \Textcite[xii]{Andrews2002Logic} uses a very appropriate term, \enquote{naive axiomatic set theory}, which we utilize.
  \end{thmenum}
\end{remark}

\begin{definition}\label{def:naive_set_theory}\mimprovised
  A particularly simple \hyperref[def:axiomatic_set_theory]{axiomatic set theory} is \term[ru=наивная теория множеств, en=naive axiomatic set theory (\Textcite[xii]{Andrews2002Logic})]{naive axiomatic set theory}, with only one existential axiom.

  \begin{thmenum}
    \thmitem{def:naive_set_theory/comprehension} The \term[en=unrestricted comprehension (\cite[240]{Mimram2020ProgramEqualsProof})]{axiom schema of unrestricted comprehension} states that any \hyperref[def:fol_formula]{first-order formula} \( \varphi \) determines a set:
    \begin{equation}\label{eq:def:naive_set_theory/comprehension}
      \qexists \syna \qforall \synx (\synx \synin \syna \syniff \syn\varphi).
    \end{equation}

    Note that \eqref{eq:def:naive_set_theory/comprehension} is a formula schema in the sense of \cref{def:fol_formula_schema}. An \hyperref[con:schemas_and_instances]{instance} of \( \syn\varphi \) may have free variables other than \( \synx \), so we must take the \hyperref[def:fol_quantifier_closure]{universal closure} of each concrete instance.

    The schema is a adaptation, in the sense of \cref{def:fol_adaptation_of_monadic_sol_theory}, of the following second-order axiom:
    \begin{equation}\label{eq:def:naive_set_theory/comprehension/sol}
      \qforall {\synp^{\syn\iota \synimplies \syn\omicron}} \qexists {\syna^{\syn\iota}} \qforall {\synx^{\syn\iota}} (\synx \synin \syna \syniff \synp \synx).
    \end{equation}
  \end{thmenum}
\end{definition}
\begin{comments}
  \item See \cref{rem:naive_set_theory} for how different authors describe this theory.
\end{comments}

\begin{definition}\label{def:set_theory_definition}
  An \hyperref[def:axiomatic_set_theory]{axiomatic set theory} \( T \) generally features no \hyperref[def:fol_signature/function]{function symbols}, making variables the only \hyperref[def:fol_term]{logical terms}.

  \hyperref[def:naive_set_theory/comprehension]{Unrestricted comprehension} and its less powerful variants like \hyperref[def:zfc/specification]{specification} allow introducing function symbols as \hyperref[def:fol_definition/function]{first-order definitions}, thus not interfering with the logical properties of the theory.

  Fix a \hyperref[def:fol_formula]{formula} \( \varphi \) whose free variables are \( x, u_1, \ldots, u_n \). The corresponding instance of unrestricted comprehension is
  \begin{equation*}
    \qforall {u_1} \ldots \qforall {u_n} \qexists a \qforall x (x \synin a \syniff \varphi).
  \end{equation*}

  With uniqueness provided by extensionality, the above acts as a proof of correctness for a definition of an \( n \)-ary \hyperref[def:fol_signature/function]{function symbol} \( A \):
  \begin{equation*}
    \qforall {u_1} \ldots \qforall {u_n} \qforall a \parens[\big]{ (A(u_1, \ldots, u_n) \syneq a) \syniff \qforall x (x \synin a \syniff \varphi)) }.
  \end{equation*}

  We can now use \( A \) to construct terms in the corresponding \hyperref[def:fol_definitional_extension]{definitional extension} \( T^+ \). Any formula and any proof in \( T^+ \) can be encoded more verbosely in \( T \).

  We call \( A \) a \term{definition by specification} and we call the variables \( u_1, \ldots, u_n \) \term{parameters}. For convenience, we will use the following metasyntax, called \term[en=set-builder notation (\cite[10]{CitkinMuravitsky2022ConsequenceRelations})]{set-builder notation}.:
  \begin{equation*}
    A(u_1, \ldots, u_n) \coloneqq \set{ x \given \varphi }.
  \end{equation*}

  Of course, if \( n = 0 \), we avoid the parameter list altogether and write
  \begin{equation*}
    A \coloneqq \set{ x \given \varphi }.
  \end{equation*}

  \begin{thmenum}
    \thmitem{def:set_theory_definition/function} Some variations of this notation are common. For example, if \( g \) is a unary function symbol, we can define \( A \) by
    \begin{equation*}
      \qforall {u_1} \ldots \qforall {u_n} \qforall a \parens[\big]{ (A(u_1, \ldots, u_n) \syneq a) \syniff \qforall x (g(x) \synin a \syniff \varphi)) }.
    \end{equation*}
    and denote it as
    \begin{equation*}
      A(u_1, \ldots, u_n) \coloneqq \set{ g(x) \given \varphi }.
    \end{equation*}

    \thmitem{def:set_theory_definition/restricted} With the axiom schema specification, which provides restricted comprehension, we must formally add an additional parameter \( u \) and substitute \( \varphi \) with \( \varphi \synwedge x \in u \), but informally we can simply write
    \begin{equation*}
      A(u_1, \ldots, u_n, u) \coloneqq \set{ x \synin u \given \varphi }.
    \end{equation*}

    If \( n = 0 \), by fixing a closed term \( \tau \) to substitute for \( u \), we can avoid the parameter list altogether and write
    \begin{equation*}
      A \coloneqq \set{ x \synin \tau \given \varphi }
    \end{equation*}

    \thmitem{def:set_theory_definition/roster} Finally, if \( \varphi = (x \syneq \tau_1) \synvee \cdots \synvee (x \syneq \tau_n) \) for closed terms \( \tau_1, \ldots, \tau_n \), we simply enumerate the terms:
    \begin{equation*}
      A \coloneqq \set{ \tau_1, \ldots, \tau_n }.
    \end{equation*}

    This is sometimes called \term[en=roster notation (\cite[10]{CitkinMuravitsky2022ConsequenceRelations})]{roster notation}.
  \end{thmenum}
\end{definition}
\begin{comments}
  \item We formalize common patterns ubiquitous in the practical use of sets. The concrete phrase \enquote{set builder notation} can be found in \cite[10]{CitkinMuravitsky2022ConsequenceRelations} and \cite[122]{Rosen2019DiscreteMathematics}.
\end{comments}

\begin{remark}\label{rem:set_existence}
  \hyperref[def:set_theory_definition]{Set-builder notation} allows us to describe sets which may not be well-defined in the underlying theory. This is convenient, although not by itself formally justified.

  Fortunately, existence proofs are often simple enough so be handled generally like in \cref{thm:zfc_existence_theorems} for \hyperref[def:zfc]{\logic{ZFC}}.

  With the constructions from this chapter, we will rarely need to concern ourselves with such existence proofs.
\end{remark}

\begin{theorem}[Russell's paradox]\label{thm:russels_paradox}
  \hyperref[def:naive_set_theory]{Naive set theory} is \hyperref[def:fol_theory/consistent]{inconsistent}. More precisely, consider the set
  \begin{equation*}
    \synR \coloneqq \set{ \synx \given \synneg (\synx \synin \synx) }.
  \end{equation*}

  Using \hyperref[def:fol_natural_deduction_system]{classical natural deduction}, we can derive both \( \synR \synin \synR \) and \( \synneg (\synR \synin \synR) \).
\end{theorem}
\begin{comments}
  \item Here \( R \) is formally defined as per \cref{def:set_theory_definition}. We do not associate it as a set in some model of set theory because, by \fullref{thm:fol_completeness}, an inconsistent theory is not satisfiable.

  \item This paradox arises from self-reference and is one of several; see \cref{ex:self_reference/russells_paradox}.
\end{comments}
\begin{proof}
  The first-order definition of \( R \) is given by
  \begin{equation*}
    \qforall \synx (\synx \synin R \syniff \synneg (\synx \synin \synx)).
  \end{equation*}

  We can derive
  \begin{equation*}
    \begin{prooftree}
      \hypo{ [\qforall \synx (\synx \synin R \syniff \synneg (\synx \synin \synx))]^u }
      \infer1[\ref{inf:def:fol_natural_deduction/forall/elim}]{ \synR \synin \synR \syniff \synneg (\synR \synin \synR) }

      \hypo{ [\synR \synin \synR]^v }
      \infer2[\ref{inf:def:propositional_natural_deduction/iff/elim_right}]{ \synneg (\synR \synin \synR) }

      \hypo{ [\synR \synin \synR]^v }
      \infer2[\ref{inf:def:propositional_natural_deduction/neg/elim}]{ \synbot }

      \infer[left label={\( v \)}]1[\ref{inf:def:propositional_natural_deduction/neg/intro}]{ \synneg (R \synless R) }
    \end{prooftree}
  \end{equation*}

  Denote the above derivation tree by \( P \). It allows us to also derive
  \begin{equation*}
    \begin{prooftree}
      \hypo{ [\qforall \synx (\synx \synin R \syniff \synneg (\synx \synin \synx))]^u }
      \infer1[\ref{inf:def:fol_natural_deduction/forall/elim}]{ \synR \synin \synR \syniff \synneg (\synR \synin \synR) }

      \hypo{ [\qforall \synx (\synx \synin R \syniff \synneg (\synx \synin \synx))]^u }
      \ellipsis {\( P \)} { \synneg (R \synless R) }

      \infer2[\ref{inf:def:propositional_natural_deduction/iff/elim_left}]{ \synR \synin \synR }
    \end{prooftree}
  \end{equation*}
\end{proof}

\begin{definition}\label{def:universal_set}\mimprovised
  As justified by \cref{def:set_theory_definition}, in \hyperref[def:naive_set_theory]{naive set theory}, we can easily define the \term[en=universal set (\cite[xvi]{Mendelson2015Logic})]{universal set}
  \begin{equation*}
    \synV \coloneqq \set{ \synx \given \syntop }.
  \end{equation*}

  Similarly to how Russell's set from \fullref{thm:russels_paradox} demonstrates the inconsistency of naive set theory, we will show in \fullref{thm:cantors_paradox} how the universal set also leads to a similar result.

  Thus, the universal set exists only as a formal concept without a well-defined interpretation.
\end{definition}

\paragraph{Power sets}

\begin{definition}\label{def:subset}\mcite[10]{Мальцев1970АлгебраическиеСистемы}
  If every member of the \hyperref[def:set]{set} \( A \) is a member of \( B \), we say that \( A \) is a \term[bg=подмножество (\cite[144]{Тагамлицки1978ИнтегралноСмятане}), ru=подмножество, en=subset (\cite[def. 2.1.3]{Rosen2019DiscreteMathematics})]{subset} of \( B \) and write \( A \subseteq B \), and similarly we say that \( B \) is a \term[bg=подмножество (\cite[144]{Тагамлицки1978ИнтегралноСмятане}), ru=надмножество, en=superset (\cite[def. 2.1.3]{Rosen2019DiscreteMathematics})]{superset} of \( A \) and write \( B \supseteq A \).

  If \( A \subseteq B \) and \( A \neq B \), we say that the subset \( A \) is \term[ru=собственное (подмножество) (\cite[34]{БелоусовТкачёв2004ДискретнаяМатематика}), en=proper (subset) (\cite[126]{Rosen2019DiscreteMathematics})]{proper} or \term[ru=строгое (подмножество) (\cite[34]{БелоусовТкачёв2004ДискретнаяМатематика}), en=strict (subset) (\cite[403]{Müller2022HandbookOfDynamicsAndProbability})]{strict} of \( B \) and write \( A \subsetneq B \).

  Formally, we can introduce the \hyperref[def:fol_signature/notation]{infix} binary \hyperref[def:fol_signature/predicate]{predicate symbol} \( {\synsubseteq} \) via the \hyperref[def:fol_definition/predicate]{first-order predicate definition}
  \begin{equation*}
    \qforall \syna \qforall \synb (\syna \synsubseteq \synb \syniff (\synx \synin \syna \synimplies \synx \synin \synb)).
  \end{equation*}
\end{definition}
\begin{comments}
  \item The subsets of a fixed set come with the structure of a \hyperref[def:boolean_algebra/complete]{complete Boolean algebra} --- see \cref{thm:boolean_algebra_of_subsets}.
\end{comments}

\begin{remark}\label{rem:strict_subset_notation}
  Based on how we denote the \hyperref[def:partially_ordered_set/strict]{strict order} of \( {\leq} \) by \( {<} \), it is natural to use \( {\subset} \) for \hyperref[def:subset]{strict subsets}.

  Unfortunately, a lot of authors use \( A \subset B \) for the nonstrict relation, so we prefer the less ambiguous \( A \subsetneq B \).

  See \cref{tab:rem:strict_subset_notation} for a usage comparison across authors.

  \begin{table}
    \centering
    \begin{tabular}{l l l}
      \toprule
      Source                                                          & Nonstrict subset    & Strict subset        \\
      \midrule
      \cite[9]{Jech2006SetTheory}                                     & \( A \subset B \)   & ---                  \\
      \cite[def. 1.3]{Rudin1976AnalysisPrinciples}                    & \( A \subset B \)   & ---                  \\
      \cite[1]{Engelking1989GeneralTopology}                          & \( A \subset B \)   & ---                  \\
      \cite[2]{Kelley1975GeneralTopology}                             & \( A \subset B \)   & ---                  \\
      \cite[91]{Halmos1960NaiveSetTheory}                             & \( A \subset B \)   & ---                  \\
      \cite[1]{Willard2004GeneralTopology}                            & \( A \subset B \)   & ---                  \\
      \cite[7]{Зорич2019АнализЧасть1}                                 & \( A \subset B \)   & ---                  \\
      \cite[59]{ИльинСадовничийСендов1985АнализТом1}                  & \( A \subset B \)   & ---                  \\
      \cite[15]{КусраевКутателадзе2005БулевозначныйАнализ}            & \( A \subset B \)   & ---                  \\
      \cite[144]{Тагамлицки1978ИнтегралноСмятане}                     & \( A \subset B \)   & ---                  \\
      \cite[3]{Jacobson1985BasicAlgebraI}                             & \( A \subset B \)   & \( A \subsetneq B \) \\
      \cite[13]{ШеньВерещагин2020НачалаТеорииМножеств}                & \( A \subset B \)   & \( A \subsetneq B \) \\
      \cite[17]{Kunen2013SetTheory}                                   & \( A \subseteq B \) & ---                  \\
      \cite[159]{Enderton1977SetTheory}                               & \( A \subseteq B \) & ---                  \\
      \cite[3]{Harzheim2005OrderedSets}                               & \( A \subseteq B \) & \( A \subset B \)    \\
      \cite[7]{Savage1998ModelsOfComputation}                         & \( A \subseteq B \) & \( A \subset B \)    \\
      \cite[135]{Kleene2002Logic}                                     & \( A \subseteq B \) & \( A \subset B \)    \\
      \cite[503]{Müller2022HandbookOfDynamicsAndProbability}          & \( A \subseteq B \) & \( A \subset B \)    \\
      \cite[def. 2.1.3]{Rosen2019DiscreteMathematics}                 & \( A \subseteq B \) & \( A \subset B \)    \\
      \cite[def. 1.1]{Александров1977ОбщаяТопология}                  & \( A \subseteq B \) & \( A \subset B \)    \\
      \cite[def. 1.1]{АлександровКолмогоров1948ТеорияМножеств}        & \( A \subseteq B \) & \( A \subset B \)    \\
      \cite[10]{Мальцев1970АлгебраическиеСистемы}                     & \( A \subseteq B \) & \( A \subset B \)    \\
      \cite[33]{БелоусовТкачёв2004ДискретнаяМатематика}               & \( A \subseteq B \) & \( A \subset B \)    \\
      \cite[def. 3.1.14]{Tao2022AnalysisI}                            & \( A \subseteq B \) & \( A \subsetneq B \) \\
      \bottomrule
    \end{tabular}
    \caption{The table for \cref{rem:strict_subset_notation} discussing how different authors denote \hyperref[def:subset]{subsets}.}\label{tab:rem:strict_subset_notation}
  \end{table}
\end{remark}

\begin{definition}\label{def:power_set}\mcite[19]{Enderton1977SetTheory}
  For every set \( A \), we define the \term[ru=булеан (\cite[34]{БелоусовТкачёв2004ДискретнаяМатематика})]{power set} \( \pow(A) \) whose members are the \hyperref[def:subset]{subsets} of \( A \).

  Formally, in accordance with \cref{def:set_theory_definition}, we can use set-builder notation to introduce a unary \hyperref[def:fol_signature/function]{function symbol} \( {\syn\pow} \):
  \begin{equation*}
    \syn\pow(a) \coloneqq \set{ x \given x \synsubseteq a }.
  \end{equation*}

  We have used the subset predicate symbol \( {\synsubseteq} \) from \cref{def:subset}.
\end{definition}
\begin{comments}
  \item In \hyperref[def:zfc]{\logic{ZFC}}, the correctness of this definition relies on the \hyperref[def:zfc/power_set]{power set axiom} --- see \cref{thm:zfc_existence_theorems/power_set}.
\end{comments}

\paragraph{Sets of sets}

\begin{remark}\label{rem:sets_of_sets}
  In a \hyperref[def:axiomatic_set_theory/atoms]{pure set theory}, the members of sets are always sets.

  When speaking about numbers, functions, symbols or abstract structures, we regard their encoding as sets is a technicality. Thus, practically, speaking about sets of numbers or functions is not fundamentally different from speaking about informal collections of numbers or functions.

  On the other hand, we often have to deal with sets whose members we also regard as sets. To avoid repetition, it is convenient to call them \enquote{collections} or \enquote{families} of sets (\enquote{\textru{совокупность}} and \enquote{\textru{семейство}} in Russian). Both terms are introduced as a synonym for \enquote{set} in
  \cite[12]{Jech2006SetTheory},
  \cite[1]{Kelley1975GeneralTopology},
  \cite[1]{Willard2004GeneralTopology},
  \cite[27]{Rudin1976AnalysisPrinciples},
  \cite[2]{Folland1999RealAnalysis} and
  \cite[5]{Зорич2019АнализЧасть1}.

  We will avoid using the phrase \enquote{family of sets} because the word \enquote{family} can also refer to indexed family as defined in \cref{def:indexed_family}. This usage is introduced, for example,
  \cite[44]{Halmos1960NaiveSetTheory},
  \cite[81]{Bourbaki2004TheoryOfSets} and
  \cite[3]{Harzheim2005OrderedSets}.
\end{remark}

\begin{definition}\label{def:empty_set}\mcite[8]{Jech2006SetTheory}
  The dual of the paradoxical \hyperref[def:universal_set]{universal set} is the \term[ru=пустое множество (\cite[6]{ШеньВерещагин2020НачалаТеорииМножеств})]{empty set} \( \varnothing \) with no members.

  Formally, in accordance with \cref{def:set_theory_definition}, we can use set-builder notation to introduce a nullary \hyperref[def:fol_signature/function]{function symbol} \( {\syn\varnothing} \):
  \begin{equation*}
    \syn\varnothing \coloneqq \set{ \synx \given \synbot }.
  \end{equation*}
\end{definition}
\begin{comments}
  \item In \hyperref[def:zfc]{\logic{ZFC}}, the correctness of this definition relies on the \hyperref[rem:zfc_formulations/empty_set]{axiom of the empty set} or an equivalent --- see \cref{thm:zfc_existence_theorems/empty_set}.

  \item This is an \hyperref[def:universal_objects/initial]{initial object} in the \hyperref[def:category_of_small_sets]{category of small sets} --- see \cref{thm:universal_sets}.
\end{comments}

\begin{definition}\label{def:subsingleton_set}\mcite[112]{UnivalentFoundationsProgram2013HoTT}
  If all elements of a set are equal, we call it a \term{subsingleton set}.
\end{definition}
\begin{comments}
  \item The term \enquote{subsingleton} is suggested as a term for \hyperref[def:mere_proposition]{mere proposition} in \cite[112]{UnivalentFoundationsProgram2013HoTT}. We adapt its usage for axiomatic set theory.

  On the other hand, \enquote{singleton} is a more popular term, given as a definition of axiomatic set theory in
  \cite[10]{Halmos1960NaiveSetTheory},
  \cite[7]{Jech2006SetTheory},
  \cite[3]{Kelley1975GeneralTopology},
  \cite[131]{КановейЛюбецкий2010СовременнаяТеорияМножеств} and
  \cite[129]{КолмогоровДрагалин2006ЛогикаВведение}.

  \item An empty set has zero elements, a singleton has one unique element, while a subsingleton may be either the empty set or a singleton. We give this more abstract definition which avoids direct references to numbers.
\end{comments}

\begin{definition}\label{def:singleton_set}\mcite[10]{Halmos1960NaiveSetTheory}
  We call a \hyperref[def:empty_set]{nonempty} \hyperref[def:singleton_set]{subsingleton set} a \term[ru=синглет (\cite[131]{КановейЛюбецкий2010СовременнаяТеорияМножеств}); синглетон (\cite[129]{КолмогоровДрагалин2006ЛогикаВведение})]{singleton}.

  Via \hyperref[def:set_theory_definition/roster]{roster notation}, we can denote the singleton containing \( a \) by \( \set{ a } \).
\end{definition}
\begin{comments}
  \item Since a singleton set is a special case of an \hyperref[def:unordered_pair]{unordered pair}, in \hyperref[def:zfc]{\logic{ZFC}}, the correctness of this definition relies on the \hyperref[def:zfc/pairing]{axiom of pairing} --- see \cref{thm:zfc_existence_theorems/singleton}.

  \item The singleton set \( \set{ A } \), where \( A \) is a \hyperref[def:small_set]{small set}, is a \hyperref[def:universal_objects/terminal]{terminal object} in the \hyperref[def:category_of_small_sets]{category of small sets} --- see \cref{thm:universal_sets}.

  \item It is sometimes convenient, especially in connection with \hyperref[ch:geometry]{geometry} or \hyperref[def:function]{set-valued maps} (e.g. when dealing with \hyperref[def:net_limit_point]{limits of nets} or \hyperref[def:subdifferentials]{subdifferentials}), to not distinguish between singleton sets and their corresponding element.
\end{comments}

\begin{definition}\label{def:unordered_pair}\mcite[9]{Halmos1960NaiveSetTheory}
  Given sets \( a \) and \( b \), via \hyperref[def:set_theory_definition/roster]{roster notation} we can define the \term[ru=неупорядоченная пара (\cite[181]{Герасимов2014Вычислимость})]{unordered pair} \( \set{ a, b } \).

  We do not require \( a \) and \( b \) to be distinct. In fact, if they are equal, we obtain a \hyperref[def:singleton_set]{singleton set}.
\end{definition}
\begin{comments}
  \item In \hyperref[def:zfc]{\logic{ZFC}}, the correctness of this definition relies on the \hyperref[def:zfc/pairing]{axiom of pairing} --- see \cref{thm:zfc_existence_theorems/unordered_pair}.

  \item We only need unordered pairs as a technical tool for defining binary unions and intersections in this sections and Kuratowski pairs in \cref{def:kuratowski_pair}. They are not by themselves an important concept.
\end{comments}

\begin{proposition}\label{thm:unordered_pair_extensionality}
  For \hyperref[def:unordered_pair]{unordered pairs} of sets the following property holds:
  \begin{equation}\label{eq:thm:unordered_pair_extensionality}
    \set{ a, b } = \set{ c, d } \T{if and only if} (a = c \T{and} b = d) \T{or} (a = d \T{and} b = c).
  \end{equation}
\end{proposition}
\begin{comments}
  \item For \hyperref[def:kuratowski_pair]{Kuratowski pairs}, we have the stronger statement \cref{thm:unordered_pair_extensionality}. For indexed families, we have \cref{thm:indexed_family_extensionality}.
\end{comments}
\begin{proof}
  \SufficiencySubProof Suppose that \( \set{ a, b } = \set{ c, d } \).

  By the \hyperref[def:naive_set_theory/extensionality]{axiom of extensionality}, both sets have the same members. We will consider several cases:
  \begin{itemize}
    \item If \( a = b \), then both \( c \) and \( d \) must be equal to their common value, hence \( a = b = c = d \). In particular, we have both
    \begin{equation*}
      a = c \T{and} b = d
    \end{equation*}
    and
    \begin{equation*}
      a = d \T{and} b = c.
    \end{equation*}

    \item If \( a \neq b \) and \( a = c \), then \( b = c \) leads to a contradiction, so it remains for \( b = d \) to hold.

    \item Similarly, if \( a \neq b \) and \( a = d \), then \( b = c \).
  \end{itemize}

  \NecessitySubProof In the converse, we also do case analysis.
  \begin{itemize}
    \item If \( a = c \) and \( b = d \), then both \( a \) and \( b \) belong to \( \set{ c, d } \), hence \( \set{ a, b } \) is a subset of \( \set{ c, d } \). We obtain the converse inclusion analogously, thus \( \set{ a, b } = \set{ c, d } \).

    \item If \( a = d \) and \( b = c \), we proceed analogously.
  \end{itemize}
\end{proof}

\paragraph{Algebra of sets}

\begin{definition}\label{def:set_union}\mcite[9]{Jech2006SetTheory}
  We will introduce two operations on sets, both called \term[ru=объединение (\cite[6]{ШеньВерещагин2020НачалаТеорииМножеств})]{union}.

  In accordance with \cref{def:set_theory_definition}, we will introduce two \hyperref[def:fol_signature/function]{function symbol} using set-builder notation:
  \begin{thmenum}
    \thmitem{def:set_union/unary} First, we introduce the \hyperref[def:fol_signature/notation]{condensed} unary function symbol \( {\synbigcup} \):
    \begin{equation*}
      \synbigcup a \coloneqq \set{ x \given \underbrace{\qexists y (y \synin a \synwedge x \synin y)}_{\qexists {y \synin a} {x \synin y}} }.
    \end{equation*}

    Informally, \( \bigcup \mscrA \) contains the common members of all members of \( \mscrA \).

    \thmitem{def:set_union/binary} Second, we introduce the infix binary function symbol \( {\syncup} \):
    \begin{equation*}
      a \syncup b \coloneqq \set{ x \given x \synin a \synwedge x \synin b }.
    \end{equation*}

    Informally, the members of \( A \cap B \) belong to both \( A \) and \( B \).
  \end{thmenum}
\end{definition}

\begin{proposition}\label{thm:binary_set_union_via_pair}
  The \hyperref[def:set_union/binary]{binary union} \( a \syncup b \) is provably equal to the \hyperref[def:set_union/unary]{unary union} of the \hyperref[def:unordered_pair]{unordered pair} \( \set{ a, b } \), i.e. we can derive
  \begin{equation*}
    a \syncup b \syneq \synbigcup \set{ a, b }.
  \end{equation*}
\end{proposition}
\begin{comments}
  \item We will prove this by constructing \hyperref[def:fol_natural_deduction_proof_tree]{natural deduction proof trees}. The proposition is much easier to state and prove in a less formal metatheory:
  \begin{equation*}
    A \cup B = \set{ x \given x \in A \T{or} x \in B } = \set{ x \given \qexists* y y \in \set{ A, B } \T{and} x \in y } = \bigcup \set{ A, B }.
  \end{equation*}
\end{comments}
\begin{proof}
  We need to derive
  \begin{align*}
    x \synin (a \syncup b)        &\syniff x \synin \synbigcup \set{ a, b }.
  \intertext{We can substitute both sides with their definitions:}
    x \synin a \synvee x \synin b &\syniff \qexists y (\underbrace{y \synin \set{ a, b }}_{y \syneq a \synvee y \syneq b} \synwedge x \synin y).
  \end{align*}

  In one direction, we have
  \begin{equation*}
    \begin{prooftree}
      \hypo{ [x \synin a \synvee x \synin b]^u }

      \infer0[\ref{inf:def:fol_natural_deduction/equality/intro}]{ a \syneq a }
      \infer1[\ref{inf:def:propositional_natural_deduction/or/intro_left}]{ a \syneq a \synvee a \syneq b }

      \hypo{ [x \synin a]^v }
      \infer2[\ref{inf:def:propositional_natural_deduction/and/intro}]{ (a \syneq a \synvee a \syneq b) \synwedge x \synin a }
      \infer1[\ref{inf:def:fol_natural_deduction/exists/intro}]{ \qexists y ((y \syneq a \synvee y \syneq b) \synwedge x \synin y) }

      \hypo{ [x \synin b]^w }
      \ellipsis {} { \synanon }

      \infer[left label={\( v, w \)}]3[\ref{inf:def:propositional_natural_deduction/or/elim}]{ \qexists y ((y \syneq a \synvee y \syneq b) \synwedge x \synin y) }
    \end{prooftree}
  \end{equation*}

  In the other direction, we have
  \begin{equation*}
    \begin{prooftree}
      \hypo{ [\qexists y ((y \syneq a \synvee y \syneq b) \synwedge x \synin y)]^u }

      \hypo{ [(y \syneq a \synvee y \syneq b) \synwedge x \synin y]^v }
      \ellipsis {\( P \)} { x \synin a \synvee x \synin b }

      \infer[left label={\( v \)}]2[\ref{inf:def:fol_natural_deduction/exists/elim}]{ x \synin a \synvee x \synin b }
    \end{prooftree}
  \end{equation*}
  where \( P \) is
  \small
  \begin{equation*}
    \begin{prooftree}[separation=1em]
      \hypo{ [(y \syneq a \synvee y \syneq b) \synwedge x \synin y]^v }
      \infer1[\ref{inf:def:propositional_natural_deduction/and/elim_left}]{ y \syneq a \synvee y \syneq b }

      \hypo{ [(y \syneq a \synvee y \syneq b) \synwedge x \synin y]^v }
      \infer1[\ref{inf:def:propositional_natural_deduction/and/elim_right}]{ x \synin y }

      \hypo{ [y \syneq a]^t }
      \infer1[\ref{inf:thm:fol_natural_deduction_equality/symm}]{ a \syneq y }
      \infer2[\ref{inf:def:fol_natural_deduction/equality/elim}]{ x \synin a }
      \infer1[\ref{inf:def:propositional_natural_deduction/or/intro_left}]{ x \synin a \synvee x \synin b }

      \hypo{ [y \syneq b]^s }
      \ellipsis {} { \synanon }

      \infer[left label={\( t, s \)}]3[\ref{inf:def:propositional_natural_deduction/or/elim}]{ x \synin a \synvee x \synin b }
    \end{prooftree}
  \end{equation*}
  \normalsize
\end{proof}

\begin{definition}\label{def:set_intersection}\mcite[14; 15]{Halmos1960NaiveSetTheory}
  Dually to unions, we will introduce two operations on sets, both called \term[ru=пересечение (\cite[6]{ШеньВерещагин2020НачалаТеорииМножеств})]{intersection}.

  In accordance with \cref{def:set_theory_definition}, we will introduce two \hyperref[def:fol_signature/function]{function symbol} using set-builder notation:
  \begin{thmenum}
    \thmitem{def:set_intersection/unary} First, we introduce the \hyperref[def:fol_signature/notation]{condensed} unary function symbol \( {\synbigcap} \):
    \begin{equation*}
      \synbigcap a \coloneqq \set{ x \given \underbrace{\qforall y (y \synin a \synimplies x \synin y)}_{\qforall {y \synin a} x \synin y} }.
    \end{equation*}

    Informally, \( \bigcap \mscrA \) contains all members of all members of \( \mscrA \) --- it \enquote{flattens} the set \( \mscrA \). Expressed more rigorously, the unary intersection is the multiplication operation of the power set monad --- see \cref{ex:def:monad/power_set}.

    \thmitem{def:set_intersection/binary} Second, we introduce the infix binary function symbol \( {\syncap} \):
    \begin{equation*}
      a \syncap b \coloneqq \set{ x \given x \synin a \synvee x \synin b }.
    \end{equation*}

    Informally, the members of \( A \cap B \) belong to either \( A \) or \( B \).
  \end{thmenum}
\end{definition}

\begin{proposition}\label{thm:binary_set_intersection_via_pair}
  The \hyperref[def:set_intersection/binary]{binary intersection} \( a \syncap b \) is provably equal to the \hyperref[def:set_intersection/unary]{unary intersection} of the \hyperref[def:unordered_pair]{unordered pair} \( \set{ a, b } \), i.e. we can derive
  \begin{equation*}
    a \syncap b \syneq \synbigcap \set{ a, b }.
  \end{equation*}
\end{proposition}
\begin{comments}
  \item We will prove this by constructing \hyperref[def:fol_natural_deduction_proof_tree]{natural deduction proof trees}. The proposition is much easier to state and prove in a less formal metatheory:
  \begin{equation*}
    A \cap B = \set{ x \given x \in A \T{and} x \in B } = \set{ x \given \qforall* y y \in \set{ A, B } \T{implies} x \in y } = \bigcap \set{ A, B }.
  \end{equation*}
\end{comments}
\begin{proof}
  We need to derive
  \begin{align*}
    x \synin (a \syncap b)        &\syniff x \synin \synbigcap \set{ a, b }.
  \intertext{We can substitute both sides with their definitions:}
    x \synin a \synwedge x \synin b &\syniff \qforall y (\underbrace{y \synin \set{ a, b }}_{y \syneq a \synvee y \syneq b} \synimplies x \synin y).
  \end{align*}

  In one direction, we have
  \begin{equation*}
    \begin{prooftree}
      \hypo{ [y \syneq a \synvee y \syneq b]^t }

      \hypo{ [x \synin a \synwedge x \synin b]^u }
      \infer1[\ref{inf:def:propositional_natural_deduction/and/elim_left}]{ x \synin a }

      \hypo{ [y \syneq a]^s }
      \infer1[\ref{inf:thm:fol_natural_deduction_equality/symm}]{ a \syneq y }

      \infer2[\ref{inf:def:fol_natural_deduction/equality/elim}]{ x \synin y }

      \hypo{ [y \syneq b]^r }
      \ellipsis {} { \synanon }

      \infer[left label={\( s, r \)}]3[\ref{inf:def:propositional_natural_deduction/or/elim}]{ x \synin y }
      \infer[left label=\( t \)]1[\ref{inf:def:propositional_natural_deduction/imp/intro}]{ (y \syneq a \synvee y \syneq b) \synimplies x \synin y }
      \infer1[\ref{inf:def:fol_natural_deduction/forall/intro}]{ \qforall \syny ((y \syneq a \synvee y \syneq b) \synimplies x \synin y) }
    \end{prooftree}
  \end{equation*}

  In the other direction, we have
  \begin{equation*}
    \begin{prooftree}
      \hypo{ [\qforall \syny ((y \syneq a \synvee y \syneq b) \synimplies x \synin y)]^u }
      \infer1[\ref{inf:def:fol_natural_deduction/forall/elim}]{ (a \syneq a \synvee a \syneq b) \synimplies x \synin a }

      \infer0[\ref{inf:def:fol_natural_deduction/equality/intro}]{ a \syneq a }
      \infer1[\ref{inf:def:propositional_natural_deduction/or/intro_left}]{ a \syneq a \synvee a \syneq b }

      \infer2[\ref{inf:def:propositional_natural_deduction/imp/elim}]{ x \synin a }

      \hypo{ [\synanon]^u }
      \ellipsis {} { x \synin b }

      \infer2[\ref{inf:def:propositional_natural_deduction/and/intro}]{ x \synin a \synwedge x \synin b }
    \end{prooftree}
  \end{equation*}
\end{proof}

\paragraph{Inductive sets}

\begin{remark}\label{rem:inductive_sets}
  We are now interested in the inductive axiom schema \eqref{eq:def:peano_arithmetic/theory/PA3}, which will allow us to define, among other things, natural numbers. Induction is an important proof technique that is discussed in detail in our proof of \cref{thm:peano_arithmetic_predecessor_existence}. There are more general forms of induction than \eqref{eq:def:peano_arithmetic/theory/PA3} like \fullref{thm:well_founded_induction} and \fullref{thm:well_founded_induction}. They do, however, require concepts which in turn depend on the existence of natural numbers within set theory. As a consequence, we cannot prove \eqref{eq:def:peano_arithmetic/theory/PA3} via \fullref{thm:well_founded_induction}.

  We will introduce the concept of inductive sets in \cref{def:inductive_set} and prove in \fullref{thm:omega_induction} that a special inductive set \hyperref[thm:smallest_inductive_set_existence]{\( \omega \)}, which will be the universe of our model of \( \BbbN \), allows performing inductive proofs. The technique that allows us to perform inductive proofs on \( \omega \) can be seen in our proof of \cref{thm:omega_is_transitive}. \Fullref{thm:omega_induction} will allow us to define natural numbers without relying on metalogical induction along the way. See our proof of \fullref{thm:omega_induction} for a description of natural number induction within set theory and \cref{rem:def:peano_arithmetic_term_model} for a further discussion of the use of natural numbers in the metatheory and in the object logic.

  We also introduce \term{recursion} in parallel as a technique for constructing objects. See \fullref{thm:omega_recursion}.
\end{remark}

\begin{definition}\label{def:ordinal_successor}\mcite[68]{Enderton1977SetTheory}
  We define the \term[ru=операция следования (\cite[1.1.11]{КолмогоровДрагалин2006ЛогикаВведение})]{successor} \( \op{sc}(A) \) of a set \( A \) as the set
  \begin{equation*}
    \op{sc}(A) \coloneqq A \cup \set{ A }.
  \end{equation*}

  The following \hyperref[con:formula_defined_predicate]{formula-defined predicate}
  \begin{equation*}\taglabel[\op{IsSucc}]{eq:def:ordinal_successor/predicate}
    \ref{eq:def:ordinal_successor/predicate}[\rho, \sigma] \coloneqq \qforall \synx \parens[\big]{ \synx \in \rho \syniff (\synx \in \sigma \synvee \synx = \sigma) },
  \end{equation*}
  which states that \( \rho \) is the successor of \( \sigma \), will be useful for us when working with \hyperref[def:inductive_set]{inductive sets}.
\end{definition}
\begin{comments}
  \item We also call it the \term{ordinal successor} operation since it is an important concept in the \hyperref[sec:ordinals]{theory of ordinals}. See \cref{rem:def:ordinal_successor} for an example of how it naturally arises. It should be distinguished from successor cardinals defined in \cref{def:successor_and_limit_cardinal}.
\end{comments}

\begin{definition}\label{def:inductive_set}\mcite[68]{Enderton1977SetTheory}
  A set is called \term{inductive} if contains the empty set and is closed under the \hyperref[def:ordinal_successor]{successor operator}.

  We introduce the following \hyperref[con:formula_defined_predicate]{formula-defined predicate}
  \begin{equation*}\taglabel[\op{IsInductive}]{eq:def:inductive_set/predicate}
    \ref{eq:def:inductive_set/predicate}[\sigma] \coloneqq
      \parens[\big]{ \qexists {\synx \in \sigma} \ref{eq:def:empty_set/predicate}[\synx] }
      \synwedge
      \parens[\big]{ \qforall {\synx \in \sigma} \qexists {\syny \in \sigma} \ref{eq:def:ordinal_successor/predicate}[\syny, \synx] }.
  \end{equation*}
\end{definition}

\begin{proposition}\label{thm:smallest_inductive_set_existence}
  There is a smallest (with respect to set inclusion) \hyperref[def:inductive_set]{inductive set}, which we denote by \( \omega \).
\end{proposition}
\begin{proof}
  We cannot directly define \( \omega \) as the intersection of all inductive sets since, in \logic{ZFC}, we want to avoid unrestricted comprehension. Fortunately, the existence of at least one inductive set is justified by the \hyperref[def:zfc/infinity]{axiom of infinity} (or by taking the entire universe in naive set theory).

  Let \( A \) be an inductive set. We can now define \( \omega \) via restricted comprehension:
  \begin{equation*}
    \omega \coloneqq \set{ x \in A \given x \T{belongs to every inductive set} }.
  \end{equation*}

  To see that \( \omega \) is itself inductive, note that \( \varnothing \in \omega \) and that if \( x \in \omega \), then it also belongs to all inductive sets and hence \( \op{sc}(x) \) also belongs to all inductive sets, proving \( \op{sc}(x) \in \omega \).
\end{proof}

\begin{theorem}[Induction via inductive sets]\label{thm:omega_induction}\mcite[thm. 6.2.17]{Hinman2005Logic}
  We can perform induction on the \hyperref[thm:smallest_inductive_set_existence]{smallest inductive set} \( \omega \). We can prove that \( \omega \) is a subset of some set \( A \) by proving the following:
  \begin{itemize}
    \item \( A \) contains \( \varnothing \).
    \item If \( A \) contains \( n \), it also contains \( \op{sc}(n) \).
  \end{itemize}

  More precisely, the following is a theorem of both naive set theory and \hyperref[def:zfc]{\logic{ZF}} (\( \omega \) satisfies \( \sigma \) and \( \synx \) plays the role of \( A \) above):
  \footnotesize
  \begin{equation}\label{eq:thm:omega_induction}
    \qexists \sigma
    \qforall \synx
    \parens[\Bigg]
      {
        \parens[\Bigg]
          {
            \parens[\big]
            {
              \underbrace{ \qexists {\syny \in \synx} \ref{eq:def:empty_set/predicate}[\syny] }_{\mathclap{\T{base case}}}
            }
            \synwedge
            \qforall \syny \parens[\big]
              {
                \overbrace
                  {
                    \underbrace{ \syny \in \synx }_{\mathclap{\substack{\T{inductive} \\ \T{hypothesis}}}}
                    \synimplies
                    \underbrace
                      {
                        \qexists {\synz \in \synx} \ref{eq:def:ordinal_successor/predicate}[\synz, \syny]
                      }_{\mathclap{\substack{\T{inductive step} \\ \T{conclusion}}}}
                  }^{\T{inductive step}}
              }
          }
        \synimplies
        \underbrace{ \ref{eq:def:subset/predicate}[\sigma, \synx] }_{\T{conclusion}}
      }
  \end{equation}
  \normalsize
\end{theorem}
\begin{comments}
  \item This is an analog of \eqref{eq:def:peano_arithmetic/theory/PA3} and is actually used in \cref{thm:omega_is_model_of_pa} to prove that \( \omega \) is a model of \hyperref[def:peano_arithmetic]{\logic{PA}}. Instead of an entire theorem schema, however, for this theorem it is sufficient to use one single formula. The more general induction principles that use theorem schemas cannot be proved without natural numbers, which are a model of \logic{PA} by virtue of this theorem.
\end{comments}
\begin{proof}
  The antecedent of (the inner formula in) \eqref{eq:thm:omega_induction} is a restatement of the formula-defined predicate \( \ref{eq:def:inductive_set/predicate}[\synx] \). The situation resembles the \hyperref[eq:def:zfc/infinity]{axiom of infinity}, but, instead of existence of an inductive set \( \synx \), it states the existence of a set \( \sigma \) such that if \( \synx \) is an inductive set, then \( \sigma \) is a subset of \( \synx \) (if we restrict \( \syny \) to range only over members of \( \sigma \), then we would obtain equality of \( \synx \) and \( \sigma \) instead). In other words, we have reduced the verification of \eqref{eq:thm:omega_induction} to showing that there exists a minimal inductive set in both naive set theory and \logic{ZF}.

  We have already proved in \cref{thm:smallest_inductive_set_existence} that our fixed model \( \mscrV = (V, I) \) of set theory has a minimal inductive set \( \omega \). Thus, for any \hyperref[def:fol_variable_assignment]{variable assignment} \( v \), the modified assignment \( v_{\sigma \mapsto \omega} \) satisfies \eqref{eq:thm:omega_induction} with the outer existential quantifier removed. Hence, \( v \) satisfies the entire formula \eqref{eq:thm:omega_induction}.

  Both the assignment \( v \) and the model \( \mscrV \) were arbitrary, therefore we can conclude that \eqref{eq:thm:omega_induction} is a theorem of both naive set theory and \logic{ZF}.
\end{proof}

\begin{definition}\label{def:transitive_set}\mcite[def. 6.1.3(i)]{Hinman2005Logic}
  We say that a set \( A \) is \term{transitive} if from \( B \in A \) it follows that \( B \subseteq A \).

  We introduce the following \hyperref[con:formula_defined_predicate]{formula-defined predicate}:
  \begin{equation*}\taglabel[\op{IsSetTransitive}]{eq:def:transitive_set/predicate}
    \ref{eq:def:transitive_set/predicate}[\sigma] \coloneqq \qforall {\synx \in \sigma} \qforall {\syny \in \synx} {\syny \in \sigma}
  \end{equation*}
\end{definition}
\begin{comments}
  \item See \cref{rem:ordinal_definition} for a discussion of the motivation and terminology of transitive sets and \cref{def:set_theory_model/transitive} for their importance.
\end{comments}

\begin{proposition}\label{thm:omega_is_transitive}
  The set \( \omega \) is transitive and every member of \( \omega \) is transitive.
\end{proposition}
\begin{comments}
  \item This proof demonstrates how \fullref{thm:omega_induction} is used. From now on, we will not be as explicit about the use of induction on \( \omega \).
\end{comments}
\begin{proof}
  \SubProof{Proof that all members of \( \omega \) are transitive} In order to demonstrate how \fullref{thm:omega_induction} works in practice, we will use inductive sets directly. Let \( T \subseteq \omega \) be the subset of all transitive members of \( \omega \). We will show that \( T \) is inductive.

  Clearly \( \varnothing \in T \)  because every member of \( \varnothing \) vacuously is a subset of \( \varnothing \).

  Now suppose that \( n \in T \) and let \( m \in \op{sc}(n) = n \cup \set{ n } \). If \( m = n \), then \( m \in \op{sc}(n) \) by definition of the successor operation. If \( m \in n \), then \( m \subseteq n \) by the inductive hypothesis and hence also \( m \subseteq n \cup \set{ n } = \op{sc}(n) \). Thus, \( \op{sc}(n) \) is also transitive.

  We have shown that \( T \) is inductive. Therefore, \( \omega \subseteq T \), i.e. every member of \( \omega \) is transitive.

  \SubProof{Proof that \( \omega \) is transitive} We will show that, for all members \( n \) of \( \omega \), we have \( n \subseteq \omega \).

  The case \( n = \varnothing \) is again trivial.

  Now suppose that \( n \subseteq \omega \) and let \( m \in \op{sc}(n) \). If \( m = n \), clearly \( m \subseteq \omega \). If \( m \in n \), then \( m \subseteq n \) and, since \( n \subseteq \omega \), we have \( m \subseteq \omega \) by transitivity of \( \subseteq \).

  Therefore, \( \omega \) is transitive.
\end{proof}

\begin{lemma}\label{thm:members_of_omega_do_not_contain_themselves}
  No element of \( \omega \) is a member of itself.
\end{lemma}
\begin{proof}
  We will again use \fullref{thm:omega_induction}.

  By definition, \( \omega \notin \omega \). Now suppose that \( n \notin n \) for some \( n \in \omega \).

  Aiming at a contradiction, suppose that \( \op{sc}(n) \in \op{sc}(n) \). The assumption that \( n = \op{sc}(n) = n \cup \set{ n } \) implies that \( n = \set{ n } \). The assumption that \( \op{sc}(n) \in n \) implies that \( n \in n \) since \( \op{sc}(n) \) is transitive by \cref{thm:omega_is_transitive}. In both cases we have \( n \in n \), which contradicts our inductive hypothesis. This contradiction shows that \( \op{sc}(n) \notin \op{sc}(n) \).

  \Fullref{thm:omega_induction} allows us to conclude that no member of \( \omega \) contains itself.
\end{proof}

\begin{theorem}\label{thm:omega_is_model_of_pa_without_operations}
  The \hyperref[thm:smallest_inductive_set_existence]{smallest inductive set} \( \omega \) satisfies the axioms \eqref{eq:def:peano_arithmetic/theory/PA1}-\eqref{eq:def:peano_arithmetic/theory/PA3} from \hyperref[def:peano_arithmetic]{Peano arithmetic} with the following interpretation:
  \begin{thmenum}
    \thmitem{thm:omega_is_model_of_pa_without_operations/zero} \hyperref[def:peano_arithmetic/signature/zero]{Zero} is interpreted as \( \varnothing \).

    \thmitem{thm:omega_is_model_of_pa_without_operations/succ} The \hyperref[def:peano_arithmetic/signature/succ]{successor} operation \( s \) is interpreted as \( \op{sc} \).
  \end{thmenum}

  We will generalize this theorem to \cref{thm:omega_is_model_of_pa} after we are able to define the arithmetic operations in \( \omega \).
\end{theorem}
\begin{proof}
  \SubProofOf{eq:def:peano_arithmetic/theory/PA1} Let \( n, m \in \omega \) and suppose that \( \op{sc}(n) = \op{sc}(m) \). If \( n = m \), there is nothing to prove.

  Suppose that \( n \neq m \). Thus, since
  \begin{equation*}
    n \cup \set{ n } = \op{sc}(n) = \op{sc}(m) = m \cup \set{ m },
  \end{equation*}
  we have both \( n \in m \) and \( m \in n \).

  \Cref{thm:omega_is_transitive} implies that \( n \) is transitive and hence \( n \in n \), which contradicts \cref{thm:members_of_omega_do_not_contain_themselves}.

  The obtained contradiction shows that \( n = m \).

  \SubProofOf{eq:def:peano_arithmetic/theory/PA2} Suppose that \( \varnothing \) has a predecessor \( n \in \omega \). Then
  \begin{equation*}
    \varnothing = \op{sc}(n) = n \cup \set{ n },
  \end{equation*}
  which implies that \( n \in \varnothing \). But this contradicts the \hyperref[def:empty_set]{definition of \( \varnothing \)}.

  Therefore, \( \varnothing \) has no predecessor.

  \SubProofOf{eq:def:peano_arithmetic/theory/PA3} It follows from \fullref{thm:omega_induction} that \eqref{eq:thm:omega_induction} is a theorem of \logic{ZF}. Let \( \mscrV = (V, I) \) be our ambient \hyperref[def:set_theory_model/standard]{standard} \hyperref[def:set_theory_model/transitive]{transitive} model of \logic{ZFC}.

  Fix any \hyperref[def:fol_variable_assignment]{variable assignment} \( v \). As in our proof of \fullref{thm:omega_induction}, we consider the modified assignment \( v_{\tau \mapsto \omega} \) that \enquote{eliminates} the outer existential quantifier in \eqref{eq:thm:omega_induction}.

  To show that \eqref{eq:thm:omega_induction} really corresponds to \eqref{eq:def:peano_arithmetic/theory/PA3} (and hence that \( \omega \) satisfies \eqref{eq:def:peano_arithmetic/theory/PA3}), fix some formula \( \varphi \) of \hyperref[def:peano_arithmetic]{Peano arithmetic} (\emph{not \logic{ZFC}!}) and suppose that \( \synx, \synz_1, \ldots, \synz_n \) are all of its free variables. Fix also some parameter values \( u_1, \ldots, u_n \in \omega \) and, as in \cref{def:set_theory_definition}, define the set
  \begin{equation*}
    A \coloneqq \set{ x \in \omega \given \Bracks{\varphi}(x, u_1, \ldots, x_n) }.
  \end{equation*}

  Since \eqref{eq:thm:omega_induction} is satisfied by \( v \), the inner formula in \eqref{eq:thm:omega_induction} (without the quantifiers over \( \tau \) and \( \sigma \)) is satisfied by \( v_{\tau \mapsto \omega, \sigma \mapsto A} \).

  Since our choice of parameters \( u_1, \ldots, u_n \) was arbitrary, we can conclude that the universal closure of \eqref{eq:def:peano_arithmetic/theory/PA3} is satisfied by \( \omega \) for every formula \( \varphi \) of \logic{PA}.
\end{proof}

\begin{remark}\label{rem:set_theory_natural_numbers_without_operations}
  Due to \cref{thm:omega_is_model_of_pa_without_operations}, we will henceforth identify the \hyperref[thm:smallest_inductive_set_existence]{smallest inductive set} \( \omega \) with the set \( \BbbN \) of \hyperref[def:natural_numbers]{natural numbers}.

  We are not yet able to add or multiply natural numbers, nor rely on their well-foundedness, however for all other purposes we are able to utilize them.

  Since the ordering in \cref{def:peano_arithmetic_ordering} is defined via addition, we must define some other ordering. Luckily, as we shall see in \fullref{sec:ordinals}, \( n < m \) corresponds to \( n \in m \). In particular, the members of \( m \) are ordered.

  As a consequence, every natural number \( n \) equals the set of all smaller natural numbers by the \hyperref[def:naive_set_theory/extensionality]{axiom of extensionality}. It is conventional to write \( \set{ 0, 1, \ldots, n } \) rather than \( \set{ m \given m \in \op{sc}(n) } \). The former notation will be fully justified in \fullref{sec:ordinals}.

  Until then, natural numbers will be useful, for example, in \cref{def:sequence}.
\end{remark}
